% Pursuing Stacks (cotes 134-2 to 134-8): archive cote 134-3, site page 20 (PDF page 21), typescript page 20.
% typescript. Draft compared: ps2.tex:909-968.
% AI-transcribed by Claude Opus 5.5 (claude-opus-5-5, adaptive thinking, effort high), Message Batches API, batch ps1 (msgbatch_01Rv5hCaNTtUt4JHx1yqoHPL), 2026-11-12; not yet checked by eye.
\scanpage{134-3/20}
\mspage{\hand{20}}

admits both types of limits,  without finitesness requirement  -  and
we certainly would expect \inserted[typed]{indeed} at least existence of finite \uline{and} infinite di-%
rect sums in M,  if objects of M were to describe arbitrary homotopy
types.  However,  in view of the special exactness properties of (Hot),
which are by no means autodual,  we will expect moreover direct sums in
M to be ``universal and disjoint'',  in Giraud's sense.  This condition,
which caracterizes to a certain extent categories which at least mildly
resemble or parallel  categories such as (Sets), (Spaces) and similar
categories,whose objects more or less express ``shapes'' -  this condition
at once rules out the majority of the most common algebraic structures,
such as rings,  groups,\xout{and the} modules over a ring or anything which
yields for M an abelian category,  etc.  If we describe an algebraic
structure species in terms of its universal realization in a category
stable under finite inverse limits,  then \deleted{a}\inserted[typed]{such} structure species can be
viewed as being defined by such a category $\underline{C}$ , and its realizations in
\inserted{any other such $C$ can be interpreted} \deleted{C} as the left-exact functors $\underline{C}\hand{\to}C$ (the universal realization of the
structure within $\underline{C}$  corresponding to the identity functor $\underline{C}\hand{\unsure{\Longrightarrow}}\underline{C}$).  In
terms of the dual category $\underline{B}$ ,  associating to every element in
Ob $\underline{B}$  =  Ob $\underline{C}$  the covariant functor $\underline{C}\hand{\to}$(Sets) it represents,  we get
a fully-faithful embedding
\[ \underline{B}\hand{\hookrightarrow} M , \]
by which $\underline{B}$ can be interpreted as the \inserted{category of the} (set-theoretic) realizations of
the given structure  which are of ``finite presentation'' in a suitable
sense (in terms of a given family of generators of $\underline{B}$ namely cogenerators
of $\underline{C}$,  considered as corresponding to the choice of ``base-sets'' for the
given structure species  -  \xout{this}\inserted[typed]{such} choice however being considered as
a convenient way merely to describe the species in concrete terms ...).
If I remember correctly,  M can be deduced from $\underline{B}$,  up to equivalence,
as being merely the category of Ind-objects of $\underline{B}$ ,  i.e. the inclusion
functor above yields an equivalence of categories
\[ \mathrm{Ind}(\underline{B}) \hand{\xrightarrow{\approx}} M , \]
which implies, I guess,  that the exactness properties of M mainly
reflect those of $\underline{B}$ .  Thus I would expect the condition we want on direct
sums in M to correspond to the same condition for finite direct sums in
$\underline{B}$ ,  not more not less.  Thus the algebraic structure species satisfying
this condition should correspond exactly to small categories $\underline{B}$,  stable
under finite direct limits,  and such that finite sums in $\underline{B}$ are disjoint
and universal.   This condition is presumably necessary,  if we want the
functor (*) from M to Hot  to be defined and to be a localization func-%
tor  -   a condition which it would be nice to understand directly in \xout{\unsure{te}}
